\section*{ВВЕДЕНИЕ}
\addcontentsline{toc}{section}{ВВЕДЕНИЕ}

В современном мире перед авторами художественных текстов стоит задача не только создавать качественный контент, но и эффективно управлять им. Написание крупного произведения требует структурирования материала, отслеживания статистики, работы с черновиками и экспорта готовых глав в распространённые форматы.

Однако существующие решения не предоставляют полноценного веб-инструмента. Ни текстовые редакторы общего назначения, ни специализированные писательские приложения не сочетают в себе иерархическую организацию текста, встроенный редактор с автосохранением, подсчёт статистики, импорт черновиков и экспорт готовых глав.

Таким образом, разработка веб-приложения, объединяющего иерархическую организацию текста, встроенный редактор с автосохранением, импорт и экспорт документов, а также ИИ-ассистента, работающего с учётом контекста книги, является актуальной задачей.

\emph{Цель настоящей работы} -- разработка веб-приложения, автоматизирующего процесс написания и структурирования книг, предоставляющего встроенного интеллектуального ассистента. Для достижения поставленной цели необходимо решить \emph{следующие задачи:}
\begin{itemize}
    \item провести анализ существующих систем и обосновать необходимость собственной разработки;
    \item спроектировать архитектуру приложения и базу данных;
    \item реализовать серверную логику на Django и клиентскую часть с редактором Quill;
    \item разработать алгоритмы автосохранения, импорта из PDF/DOCX и экспорта в DOCX;
    \item интегрировать ИИ-ассистента через OpenRouter API;
    \item провести функциональное тестирование.
\end{itemize}

\emph{Структура и объем работы.} Отчет состоит из введения, 4 разделов основной части, заключения, списка использованных источников, 2 приложений. Текст выпускной квалификационной работы равен \formbytotal{lastpage}{страниц}{е}{ам}{ам}.

\emph{Во введении} сформулирована цель работы, поставлены задачи разработки, описана структура работы, приведено краткое содержание каждого из разделов.

\emph{В первом разделе} на стадии анализа предметной области приводится характеристика процесса работы автора над книгой, выявляются проблемы традиционного подхода, обосновывается необходимость автоматизации и анализируются существующие аналоги.

\emph{Во втором разделе} на стадии технического задания формулируются требования к функциональности, интерфейсу, надёжности, информационной безопасности и совместимости.

\emph{В третьем разделе} на стадии технического проектирования представлены архитектура программной системы, обоснование выбора технологий, проект базы данных, а также проектирование маршрутов API и пользовательского интерфейса.

\emph{В четвертом разделе} приводится спецификация программных компонентов и результаты функционального тестирования в виде тестовых сценариев.

В заключении излагаются основные результаты работы, полученные в ходе разработки.

В приложении А представлен графический материал.
В приложении Б представлены фрагменты исходного кода. 
